\section{Conclusion}
\label{sec:conclusion}

We present the first calibration study of AI expert panel scores against human expert judgment in a creative domain. Across 10 puzzle hunt puzzles scored by an AI panel (using v1 and v2 expert persona profiles) and five experienced human constructors on an identical six-dimension rubric, we find a systematic and interpretable reliability-by-dimension pattern: AI panels are strong proxies for human judgment on technically verifiable quality dimensions (Clarity, $r = 0.88$; Solvability, $r = 0.84$; Confirmation Quality, $r = 0.86$) and weaker proxies on hedonically evaluated dimensions (Fun, $r = 0.52$; Elegance, $r = 0.68$). AI panels systematically underestimate Fun ($\Delta = -0.83$, $p < 0.01$) and Elegance ($\Delta = -0.58$, $p < 0.05$), with no significant bias on technical dimensions.

The profile version analysis provides a theoretically significant result: richer v2 expert personas substantially close the Elegance calibration gap (correlation improves from 0.49 to 0.68; bias narrows by 48\%) but leave the Fun calibration gap essentially unchanged (correlation 0.51 to 0.52; bias $-0.91$ to $-0.88$). This dissociation confirms that Elegance has a structural component accessible to persona-based evaluation, while Fun has a hedonic component that profile richness cannot substitute for.

\textbf{Practical guidance for practitioners.} Use AI panel scores on Clarity, Solvability, and Confirmation Quality as-is; no correction is warranted. For Fun and Elegance, apply mean-bias corrections: add $+0.83$ to AI Fun scores and $+0.58$ to AI Elegance scores when comparing against the 22/30 threshold (these are mean-shift corrections; per-puzzle slope correction requires larger $N$). Lower the AI-only pass threshold from 22/30 to approximately 20/30, or treat puzzles in the 20--22/30 range as requiring mandatory human review. For final decisions where Fun and Elegance are primary criteria, supplement AI review with human assessment.

\textbf{A generalizable hypothesis.} Our results suggest a hypothesis: AI evaluation reliability tracks criterion verifiability. Dimensions assessable by inspecting the artifact's structure are reliably AI-evaluable; dimensions requiring hedonic experience are not. This pattern is consistent with theoretical accounts of the explanatory gap between structural description and hedonic experience, and we predict it will replicate across creative evaluation domains---academic paper review, game design, interactive fiction---wherever rubrics mix technically verifiable and hedonically evaluated criteria. Multi-domain replication is needed to elevate this from a well-supported domain-specific finding to an established general principle.

\textbf{Open questions.} Three questions most directly extend this work. First, does the reliability-by-dimension pattern hold for larger AI panels (N=10+ reviewers)? Larger panels would have smaller variance in AI mean scores and might show higher correlation with human means, particularly for hedonic dimensions. Second, do the corrective scaling factors generalize across puzzle hunt communities with different aesthetic traditions? The current calibration is community-specific; cross-community replication would establish whether the bias pattern is domain-general or community-contingent. Third, does the structural/hedonic distinction predict AI evaluation reliability in other creative domains? A comparative calibration study across puzzle design, game level design, and interactive narrative would test the generalizability of the principle beyond the domain where it was first documented.

\textbf{Research artifacts.} The full rubric definition, puzzle materials (with solutions withheld), AI panel score matrices (v1 and v2), human reviewer score matrices, analysis scripts, and corrective scaling factor derivation are available at [URL]. The calibration methodology is documented as a reusable protocol applicable to any AI review system operating on a structured rubric.
